\scanpage{107}
% archive foliation: 106
\uline{Proposition 2} Les $a_\alpha$ $(1\leq\alpha\leq3)$ sont caractérisés par
la condition que $\wp^{\unsure{2|3}}+a_1\wp^2+a_2\wp+a_3\mid J=0$.

C'est nécessaire, puisque $\wp'\mid J=0$. Mais
d'après le cor. à la prop. précédente, les
fonctions $1,\wp,\wp^2$ restreintes à $J$ sont lin\textsuperscript{t}
indépendantes, donc la relation précédente
sur les $a_i$ (qui s'écrit
\[
(-\wp^3)|J=a_1(\wp^2|J)+a_2(\wp|J)+a_3(1_E|J)
\]
)
caractérise déjà les $a_i$, cqfd.

\uline{Corollaire 1} Les $\inserted{(-1)^\alpha}a_\alpha$ $(1\leq\alpha\leq3)$ sont les fonctions
sym. élémentaires en les $\wp(e_i)$ $(i\in J)$
i.e. on a
\[
(29)\qquad\wp^3+a_1\wp^2+a_2\wp\unsure{-|+}a_3
=\prod_{i\in J}(\wp-\deleted{\unsure{\wp}}\wp(e_i)).
\]
\deleted{\unsure{NB}}
\uline{Corollaire 2} Pour que l'on ait $a_1=0$, il f. et s.
que
\[
\operatorname{Tr}_{J/k}(\wp|J)=0\qquad\text{i.e. }\sum\wp(e_i)=0.
\]
% margin position: left of Corollaire 2
\margin{\figure{\begin{tikzpicture}\draw (0,0)--(0,1.5);\end{tikzpicture}}}
Si car. $k\ne3$, les $f\in\mathcal F_2$ qui satisfont cette
condition forment un supplémentaire
$W_E$ de $\mathcal F_1$ dans $\mathcal F_2$
\[
(30)\qquad\mathcal F_2\simeq\underbrace{\mathcal F_1}_{k.1_E}\oplus\underbrace{W_E}_{\simeq t_E^{\otimes2}}
\]
les \deleted{\unsure{sections}}\inserted{éléments} de $W_E$ sont \inserted{donc} les fonctions de Weierstrass,
